\section{Ruling Out the Obvious Objections}
\label{sec:analysis}

A negative result of this shape invites four specific rebuttals: the solver is
memorizing, its strength is generic to any non-learning method, GRPO's score is
a reward-hacking artifact, and GRPO is simply mistuned. Each is testable; none
survives. We then report the one run we excluded, ask whether our GRPO is
simply under-trained, and state the solver's cost.

\subsection{``The solver is just retrieving gold answers''}

If the solver's accuracy came from copying gold programs out of its retrieval
library, it would be a memorization artifact rather than a baseline. We
decompose it into three variants: \textbf{A} uses only the hand-written rule
path; \textbf{B} uses only retrieval; \textbf{C} is the deployed configuration
(rules first, retrieval as fallback).

\begin{table}[t]
\centering
\small
\caption{Solver ablation (dev, mean of 3 seeds). Retrieval never fires: its
marginal contribution $C-A$ is $+0.00$ points at all ten Config~1 budgets.}
\label{tab:ablation}
\begin{tabular}{rrrr r}
\toprule
$m$ & A: rules & B: retrieval & C: deployed & $C-A$ \\
\midrule
\multicolumn{5}{l}{\textit{Config~1 (all 10 budgets identical)}} \\
$0 \ldots 2446$ & 72.00 & 0.00 & 72.00 & $+0.00$ \\
\midrule
\multicolumn{5}{l}{\textit{Config~2}} \\
128 & 36.00 & 0.00 & 35.33 & $-0.67$ \\
256 & 35.67 & 0.00 & 35.33 & $-0.33$ \\
\bottomrule
\end{tabular}
\end{table}

Table~\ref{tab:ablation} shows the retrieval-only variant scores $0.00\%$
everywhere: it never produces an accepted program. Across all ten Config~1
budgets and three seeds---thirty measurements---$A$, $C$ and their difference
are constant at $72.00\%$, $72.00\%$ and $+0.00$ points. In Config~2 the
deployed variant differs from the rule-only path by $-0.33$ to $-0.67$
points---one to two items on the 200-item dev subset, i.e.\ within noise---so
retrieval contributes nothing measurable on either configuration. The solver's
ability is therefore entirely symbolic. This
also explains the flat dose curve: with retrieval never firing, the
demonstration budget has no channel through which to affect the solver at all.

\subsection{``Any non-learning baseline would look that strong''}

A second objection is that the solver's advantage demonstrates nothing about
symbolic competence: on a task whose answers are executable programs, perhaps
any method that skips training lands near the same place. The natural test is
the other obvious non-learning option: prompt the base model directly to emit a
program, with the same $m$ demonstrations used as few-shot examples. This is
Program-of-Thought / PAL prompting~\cite{pot2022,pal2023}---no gradient steps
at all.

\begin{table}[t]
\centering
\small
\setlength{\tabcolsep}{4pt}
\caption{Frozen-test accuracy (\%) of the two non-learning options. Prompting
the base model directly is by far the weakest arm in both configurations, so
the solver's advantage does not come free with skipping training.}
\label{tab:pot}
\begin{tabular}{lrr}
\toprule
& Config 1 & Config 2 \\
\midrule
solver (symbolic rules)          & 64.54 & 32.08--32.83 \\
PoT / PAL prompting (no training) & 9.93--16.67 & 0.75--1.36 \\
\midrule
GRPO (RL)                        & 58.51--63.12 & 16.11--30.57 \\
\bottomrule
\end{tabular}
\end{table}

Table~\ref{tab:pot} gives the result over the eleven groups of
Table~\ref{tab:main}. Few-shot prompting reaches $13.30\%$ on average in
Config~1 and $1.00\%$ in Config~2---the weakest arm in the entire study, far
below both the solver and GRPO. Two things follow. First, the solver's
strength is specific rather than generic: the most obvious alternative
non-learning baseline is dramatically worse, so what the solver contributes is
genuine symbolic competence rather than a free advantage of skipping training. Second, RLVR is doing substantial work---GRPO improves over
prompting the same base model by 44 to 50 points across the Config~1
groups---which is precisely why the comparison matters: a method that clearly
helps can still be the wrong way to spend the budget.

\subsection{``GRPO's score is inflated by verifier hacking''}

If GRPO were exploiting the reward checker---emitting constant programs,
inlining the gold value, or collapsing onto one template---its accuracy would
be an artifact and the comparison meaningless. We audit reward-dump records across the 14 GRPO runs whose reward dumps were retained, spanning all three backbones.

\begin{table}[t]
\centering
\small
\setlength{\tabcolsep}{4pt}
\caption{Verifier-hacking audit over the 14 GRPO runs whose reward dumps were
retained. H1: constant program; H4: no-op arithmetic; H3: share of positive
samples under the single most frequent code skeleton.}
\label{tab:hacking}
\begin{tabular}{lrrrrr}
\toprule
& samples & positive & H1 & H4 & H3 \\
\midrule
Config~1 (6 runs)      & 87{,}872 & 47{,}616 & 0 & 4  & $\leq 0.9\%$ \\
Config~2 (6 runs)      & 37{,}933 & 6{,}030  & 0 & 23 & $\leq 1.5\%$ \\
Config~3, 7B (2 runs)  & 4{,}672  & 3{,}871  & 0 & 0  & $\leq 2.1\%$ \\
\bottomrule
\end{tabular}
\end{table}

Table~\ref{tab:hacking} reports the result. Constant-program shortcuts occur
\textbf{zero} times in 57{,}517 positive-reward samples---the verifier's
anti-constant guard holds. No-op arithmetic appears 27 times ($0.05\%$).
Template collapse is absent: the most frequent code skeleton accounts for at
most $2.1\%$ of positive samples, with skeleton Gini coefficients between
$0.06$ and $0.45$, indicating a broadly distributed solution space.
Gold-literal matches are rare ($0.19\%$) and expected on these tasks, where
gold answers are computed from context numbers.

GRPO's $16.1$--$80.9\%$ frozen-test accuracies across the three backbones
therefore reflect genuine program synthesis. That is precisely what makes the comparison meaningful: a
policy that is really solving the task still loses to a solver that never
learned.

\subsection{``You did not tune GRPO properly''}

A fourth objection is that the negative result holds only at our particular
hyperparameters. We sweep the knob most likely to matter---the KL coefficient
$\beta$---at Config~1's higher operating budget, in its smallest-margin seed
($m{=}64$, seed~2), and evaluate every resulting policy on both dev and the
frozen test set.

\begin{table}[t]
\centering
\small
\setlength{\tabcolsep}{3.5pt}
\caption{GRPO hyperparameter sweep, Config~1 $m{=}64$, seed 2 (solver: 72.00
dev / 64.54 test). $\beta{=}0$ is the only setting whose $\Delta$ turns
positive---and only on dev.}
\label{tab:sweep}
\begin{tabular}{llrrrr}
\toprule
$\beta$ & temp & dev & test & $\Delta_{\text{dev}}$ & $\Delta_{\text{test}}$ \\
\midrule
0.04 & 1.0 & 68.00 & 61.70 & $-4.00$ & $-2.84$ \\
0.0  & 1.0 & 73.50 & 62.41 & $\mathbf{+1.50}$ & $\mathbf{-2.13}$ \\
0.2  & 1.0 & 58.00 & 51.42 & $-14.00$ & $-13.12$ \\
\bottomrule
\end{tabular}
\end{table}

Table~\ref{tab:sweep} contains the sharpest single observation in this paper.
Removing the KL penalty entirely ($\beta{=}0$) does lift GRPO above the solver
\emph{on dev}, by $1.50$ points---the only positive $\Delta$ any setting
produced. On the frozen test set the same policy scores $62.41\%$ against the
solver's $64.54\%$: $\Delta = -2.13$ points. The reversal is a selection
artifact, not an effect.

This is one of three independent instances of the same phenomenon
(Section~\ref{sec:results}): Config~1 at $m{=}2$ gave $+0.50$ on dev against
$-3.90$ on the frozen test, and the 7B $m{=}2$ seed-2 group gave $+1.50$
against $-1.77$. In all three, dev favored RL and the frozen test did not. Had we followed the common practice of reporting dev numbers after a
hyperparameter search, we would have published the opposite conclusion---from
the same runs.

\subsection{One excluded run}

One group (Config 2, $m=128$, seed 1) experienced training collapse, ending with a zero positive-reward rate. We exclude this run as an infrastructure/training failure rather than an algorithmic property; notably, including it ($\Delta = -19.73$) would further strengthen our negative finding.

\begin{table}[t]
\centering
\small
\caption{Positive-reward rate across all 18 GRPO runs. One run, on the 1B
backbone at the scarcest budget, ends at zero and is excluded.}
\label{tab:stability}
\begin{tabular}{lrrr}
\toprule
Backbone & budget & runs & positive-reward rate \\
\midrule
Qwen2.5-0.5B & 2, 64   & 6 & 0.522--0.560 \\
Llama-3.2-1B & 128     & 3 & 0.081, 0.092, \textbf{0.000} \\
Llama-3.2-1B & 256     & 3 & 0.194--0.239 \\
Qwen2.5-7B   & 2, 64   & 6 & 0.772--0.846 \\
\bottomrule
\end{tabular}
\end{table}

Table~\ref{tab:stability} places it in context: the collapse is confined to the
weakest backbone at the scarcest budget, and every other run in the study
sustains a positive reward rate throughout training.

\subsection{Is our GRPO simply under-trained?}

At 7B, GRPO ranks third of five arms on seed-averaged accuracy, behind
continued SFT and RS-SFT. Since
beating SFT is RLVR's core claim, a reader should ask whether our GRPO is
broken rather than our conclusion interesting. Four pieces of evidence say it
is not.

First, \emph{it works where it should}: at 0.5B, GRPO beats every learning arm
by wide margins---$62.41\%$ against warm-start's $15.60\%$ and continued SFT's
$24.47\%$ at $m{=}2$. A broken implementation would not produce a
35-point gain in one regime and a 3-point deficit in another.

Second, \emph{the pre-registered budget was consumed}: every GRPO run reached
its 5700-second cap with 2240--17920 verifier calls and 131k--947k generated
tokens. Two qualifications are needed here, both of which the conference
version of this paper omitted.

The first concerns what ``matched compute'' means. GRPO receives the same wall
clock as continued SFT, which is the arm attaining $\max\mathcal{B}$ in all six
7B groups, and vastly more than the solver, which attains it in the other
eleven and runs on CPU in seconds. But GRPO does \emph{not} exceed every arm on
every axis: RS-SFT generates three to seven times more tokens than GRPO in
Config~1 (Table~\ref{tab:compute}). RS-SFT never attains $\max\mathcal{B}$ in
any of the 17 groups, so this does not enter $\Delta$---but the blanket claim
that GRPO received at least as much compute as the arms it is compared with is
true only on the wall-clock axis and only against the arms that actually bind.

% 由 analysis/gen_extended_stats.py 生成，勿手改
\begin{table}[t]
\centering
\small
\setlength{\tabcolsep}{3pt}
\caption{Compute on four axes, ranges over the three seeds. Wall clock is the pre-registered matching axis; the other three are reported so that the matching is not taken on trust. RS-SFT, not GRPO, is the most expensive arm on generated tokens and verifier calls.}
\label{tab:compute}
\begin{tabular}{llrrrrrrr}
\toprule
& & \multicolumn{3}{c}{wall clock (s)} & \multicolumn{2}{c}{verifier calls} & \multicolumn{2}{c}{generated tokens} \\
\cmidrule(lr){3-5}\cmidrule(lr){6-7}\cmidrule(lr){8-9}
cfg & $m$ & c-SFT & GRPO & RS-SFT & GRPO & RS-SFT & GRPO & RS-SFT \\
\midrule
C1 & 2 & 2579--3196 & 5701--5703 & 5715--5732 & 15488--17920 & 21824--24576 & 722k--798k & 5441k--6136k \\
C1 & 64 & 5703--5707 & 5711--5728 & 4561--5815 & 12160--13184 & 13184--15872 & 908k--947k & 2906k--3301k \\
C2 & 128 & 5704--5715 & 5709--5751 & 1549--5905 & 6272--7040 & 7040--9728 & 296k--323k & 1103k--1298k \\
C2 & 256 & 5703--5714 & 5720--5755 & 5822--6038 & 7104--7872 & 6144--6656 & 299k--335k & 602k--762k \\
C3 & 2 & 5702--5706 & 5712--5761 & --- & 2304--2880 & --- & 131k--150k & --- \\
C3 & 64 & 5721--5748 & 5750--5864 & --- & 2240--2496 & --- & 167k--181k & --- \\
\bottomrule
\end{tabular}
\end{table}


The second qualification is more serious, and it is the strongest limitation in
this study. Each prompt is sampled 8 times, so the verifier-call counts above
imply that GRPO visited 280--2240 distinct prompts. Against pools of 2846 and
4035 items, \textbf{no GRPO run in this study completed even one epoch}:

\begin{center}
\begin{tabular}{llrrr}
\toprule
config & $m$ & verifier calls & distinct prompts & epochs over pool \\
\midrule
Config~1 (0.5B) & 2   & 15488--17920 & 1936--2240 & 0.68--0.79 \\
Config~1 (0.5B) & 64  & 12160--13184 & 1520--1648 & 0.53--0.58 \\
Config~2 (1B)   & 128 & 6272--7040   & 784--880   & 0.19--0.22 \\
Config~2 (1B)   & 256 & 7104--7872   & 888--984   & 0.22--0.24 \\
Config~3 (7B)   & 2   & 2304--2880   & 288--360   & 0.10--0.13 \\
Config~3 (7B)   & 64  & 2240--2496   & 280--312   & 0.10--0.11 \\
\bottomrule
\end{tabular}
\end{center}

Coverage is worst exactly where the paper's generalization claim lives: at 7B,
GRPO saw roughly a tenth of the available prompts. Equal wall clock at a fixed
per-token cost buys an order of magnitude less data at 7B than at 0.5B, and our
pre-registration fixed the seconds rather than the data. We do not think this
explains the result---the argument continues below---but a reader who believes
that RLVR needs multiple epochs to pay off has not been answered by this study,
and we do not claim otherwise.

Third, \emph{the policy is genuinely solving}: the audit of
Table~\ref{tab:hacking} covers the 7B runs directly and finds zero
constant-program shortcuts, zero no-op arithmetic, and a top-skeleton share of
$2.1\%$ over 3871 positive samples.

% Generated by analysis/gen_extended_stats.py. Do not edit by hand.
\begin{table}[t]
\centering
\small
\setlength{\tabcolsep}{4pt}
\caption{Zero-advantage groups across all 18 GRPO runs, read from each run's recorded metadata. A prompt group whose 8 samples all receive the same reward has zero group-relative advantage and contributes no gradient. The 7B runs have the highest positive-reward rate in the study and the fewest informative prompts.}
\label{tab:zeroadv}
\begin{tabular}{llrrrrrrr}
\toprule
backbone & $m$/seed & steps & groups & all-0 & all-1 & zero-var & informative & pos.\ rate \\
\midrule
Qwen2.5-0.5B & 2/0 & 242 & 1936 & 565 & 652 & 62.9\% & \textbf{719} & 0.527 \\
 & 2/1 & 280 & 2240 & 625 & 871 & 66.8\% & \textbf{744} & 0.560 \\
 & 2/2 & 258 & 2064 & 572 & 748 & 64.0\% & \textbf{744} & 0.551 \\
 & 64/0 & 190 & 1520 & 381 & 387 & 50.5\% & \textbf{752} & 0.522 \\
 & 64/1 & 197 & 1576 & 348 & 427 & 49.2\% & \textbf{801} & 0.544 \\
 & 64/2 & 206 & 1648 & 386 & 441 & 50.2\% & \textbf{821} & 0.540 \\
Llama-3.2-1B & 128/0 & 98 & 784 & 544 & 1 & 69.5\% & \textbf{239} & 0.081 \\
 & 128/1$^{\dagger}$ & 110 & 880 & 880 & 0 & 100.0\% & \textbf{0} & 0.000 \\
 & 128/2 & 102 & 816 & 549 & 2 & 67.5\% & \textbf{265} & 0.092 \\
 & 256/0 & 111 & 888 & 463 & 16 & 53.9\% & \textbf{409} & 0.196 \\
 & 256/1 & 111 & 888 & 461 & 12 & 53.3\% & \textbf{415} & 0.194 \\
 & 256/2 & 123 & 984 & 460 & 36 & 50.4\% & \textbf{488} & 0.239 \\
Qwen2.5-7B & 2/0 & 38 & 304 & 29 & 207 & 77.6\% & \textbf{68} & 0.812 \\
 & 2/1 & 36 & 288 & 42 & 185 & 78.8\% & \textbf{61} & 0.772 \\
 & 2/2 & 45 & 360 & 43 & 267 & 86.1\% & \textbf{50} & 0.823 \\
 & 64/0 & 35 & 280 & 18 & 180 & 70.7\% & \textbf{82} & 0.846 \\
 & 64/1 & 39 & 312 & 24 & 206 & 73.7\% & \textbf{82} & 0.837 \\
 & 64/2 & 35 & 280 & 21 & 194 & 76.8\% & \textbf{65} & 0.826 \\
\bottomrule
\end{tabular}
\\[2pt]\footnotesize $^{\dagger}$The excluded run: every group all-zero.
\end{table}


The conference version of this paper continued the third argument by noting
that the 7B positive-reward rate, $0.772$--$0.846$, is the highest of the three
backbones, and read that as evidence that the 7B policy was training
successfully. \textbf{That inference is backwards, and we retract it.} GRPO's
advantage is computed within a group of samples drawn from the same prompt; if
all eight samples receive the same reward the group's advantage is zero and it
contributes no gradient. A high positive-reward rate therefore implies
\emph{more} degenerate groups, not more learning.

Table~\ref{tab:zeroadv} makes the size of the effect visible. At 7B, between
$70.7\%$ and $86.1\%$ of prompt groups are all-correct or all-wrong, leaving
\textbf{50 to 82 informative prompts for an entire run}. At 0.5B the same
quantity is 719 to 821. The 7B runs are thus training on roughly an order of
magnitude less gradient signal than the 0.5B runs that this study reads as its
positive control. The excluded Config~2 run appears here in its proper form:
$880$ of $880$ groups all-zero, zero informative prompts.

The honest position is that the third argument survives only in its first half.
The 7B policy is solving the task rather than gaming the verifier, and that
much the audit establishes. It is \emph{not} established that the 7B GRPO arm
received enough gradient signal to represent what GRPO can do at that scale.

We subsequently tested that directly rather than leaving it as a caveat, and the
answer is that the 7B arm was indeed under-trained: at ten times the coverage
the smallest-margin 7B cell reverses from $-1.06$ to $+6.74$
(Appendix~\ref{app:followup}). Two things about that result belong here. It is a
\emph{budget} effect and not a degeneracy effect---a matched-budget run that
raises informative prompts from 82 to 120 moves $\Delta$ to $+1.06$ with an
interval of $[-2.13,+4.26]$, which is the same place as the original. And the
degeneracy does not improve with training: across $10.6\times$ the steps the
zero-advantage fraction \emph{rose} from $73.7\%$ to $86.0\%$, so the win came
from brute force rather than from better use of the signal. Section~\ref{sec:scale}
states what this costs the 7B configuration.

Fourth, \emph{tuning does not rescue it}: the sweep in
Table~\ref{tab:sweep} covers the KL coefficient. The
single setting that lifts GRPO above the best baseline does so only on dev, and
reverses on the frozen test.

The more economical reading is that a stronger backbone helps plain imitation
at least as much as it helps RL. At 7B, continued SFT on $m$ demonstrations is
already close to what the model can do on this task; the room RL has to add
value has narrowed, while its cost has not.

\subsection{``You left out the obvious inference-time baseline''}
\label{subsec:bon}

A fifth objection, which we did not anticipate and which reviewers of the
conference version raised independently, is that the control group omits the
cheapest non-learning option of all: sample $n$ completions from the warm start
and keep the one the verifier accepts. If verifier-filtered best-of-$n$ were
strong, it would be a far more natural control arm than a hand-written solver.

It is not available on these tasks, for a structural reason. Our verifier
decides correctness by executing the emitted program and comparing its output
to the \emph{gold} answer within a 1\% relative tolerance. At training time
that is legitimate---the gold answer is part of the training item. At inference
time it is not: selecting among $n$ candidates with this verifier requires
knowing the answer. Verifier-filtered best-of-$n$ on these tasks is therefore
identical to pass@$n$, an oracle upper bound rather than a deployable arm. The
deployable relative is majority voting over the $n$ samples, which needs no gold.

We measured both from the warm start at $n{=}16$ on dev:

\begin{center}
\begin{tabular}{llrrrr}
\toprule
config & $m$ & pass@1 & pass@16 (oracle) & maj@16 (deployable) & solver \\
\midrule
Config~1 & 2   & 9.83  & 48.83 & 35.33 & 72.00 \\
Config~1 & 64  & 13.17 & 61.17 & 45.17 & 72.00 \\
Config~2 & 128 & 5.50  & 34.50 & 17.17 & 35.33 \\
Config~2 & 256 & 16.00 & 51.67 & 30.67 & 35.33 \\
\bottomrule
\end{tabular}
\end{center}

At all four operating points the deployable arm (maj@16) is below the solver,
by $15$ to $37$ points. At three of the four, so is the \emph{oracle} upper
bound: even an omniscient selector over 16 samples from the warm start would
not reach the solver at Config~1 $m{=}2$, Config~1 $m{=}64$, or Config~2
$m{=}128$. Only Config~2 $m{=}256$ has oracle headroom above the solver, and
there the deployable version still falls $4.7$ points short.

Two caveats. These are dev numbers on the 200-item selection subset, not frozen
test; and this arm was not part of the pre-registered control group, so it
enters as supporting evidence rather than into $\Delta$. Appendix~\ref{app:pot}
gives the full budget sweep, in which the sampling arms do overtake the solver
past $m{=}256$---the same crossing the dose curves show for every learning arm.

\subsection{How the solver was built, and what it cost}
\label{subsec:solvercost}

Because the solver attains $\max\mathcal{B}$ in 11 of 17 groups, its
construction is the single largest discretionary choice in this study, and a
reader is entitled to the procedure rather than an assurance. The pre-registered
protocol constrained the solver in one direction only---it forbade weakening it
if GRPO won---which does nothing to prevent the opposite failure. What actually
constrains it is where it was allowed to look.

The solver was developed against a dev split carved from training data (for
Config~2, the self-consistent subset of \texttt{train} at \texttt{idx \% 4 == 0},
$n{=}1013$), with leave-one-out retrieval, under a fast in-process replica of
the verifier's semantics. The frozen test set was evaluated \emph{once}, at the
end, with the real verifier. Development proceeded over numbered rounds whose
dev accuracy we logged as it moved---$0.243 \rightarrow 0.294 \rightarrow 0.317
\rightarrow \ldots$ for Config~2---together with the failure mode each round
targeted. Those logs are published verbatim
(Appendix~\ref{app:reproduce}); they are the audit trail for this arm, and they
are not flattering in places.

The cost of that dev fitting is visible and we should name it. The solver scores
$72.00\%$ on Config~1 dev and $64.54\%$ on Config~1 frozen test: a drop of
$7.46$ points, larger than 15 of the 17 effect sizes this paper reports. That
gap is the solver's own generalization penalty from having been iterated
against dev, and it is already fully paid inside every number in
Table~\ref{tab:main}, because $\Delta$ is computed on test. A reader who
suspects the solver of being dev-overfitted is right, and the frozen-test
protocol has already charged it for that.

The engineering cost was roughly 440 lines of task-specific adaptation code on
top of a reusable symbolic execution core, developed over the logged rounds,
after which inference runs in seconds on CPU. We do not claim this is cheap in
human terms, and ``simpler'' in this paper's title refers to the method's
inference-time and data requirements, not to the labour of writing it. Whether
a practitioner should build such an arm is a different question from whether a
paper claiming an RLVR gain should have to clear it.

\paragraph{Does the solver transfer?}
The sharpest form of the objection is that 440 dataset-specific lines are not
``simpler'' in any useful sense, because they buy nothing outside the dataset
they were written for. This is testable within our own study: we have two
tasks and two independently developed solvers, so each can be run, unmodified,
on the other's frozen test set under the same clean-subset protocol.

% 由 analysis/solver_transfer_cross.py 生成，勿手改
\begin{table}[t]
\centering
\small
\caption{Solver transferability. Each solver is run unmodified on the other task's frozen test set, under the same clean-subset protocol. The narrative solver transfers and still beats all three learning arms of Config~1 at $m{=}2$; the table solver does not transfer at all, because it cannot parse a context that is not a JSON table and abstains on every item.}
\label{tab:transfer}
\begin{tabular}{lrrr}
\toprule
solver & home task & other task & retained \\
\midrule
built on CodeTAT-QA & 64.54 & 0.00 & 0\% \\
built on CodeFinQA & 33.43 & 36.17 & 108\% \\
\bottomrule
\end{tabular}
\end{table}


Table~\ref{tab:transfer} shows that the answer depends on direction, and that
the objection is half right. The CodeTAT-QA solver does not transfer at all: it
parses a JSON table out of the context and emits cell-indexed programs, and on
CodeFinQA's narrative contexts it finds no table and abstains on all 664 items.
That is an architectural incompatibility rather than a graceful degradation.

The CodeFinQA solver transfers. Run unmodified on CodeTAT-QA it scores
$36.17\%$, slightly \emph{above} its $33.43\%$ on its home task. The comparison
that matters is with Table~\ref{tab:main}: at Config~1's $m{=}2$ operating
point the three learning arms score $14.54$--$26.60\%$ ($W_m$),
$18.09$--$24.47\%$ (c-SFT) and $27.66$--$30.50\%$ (RS-SFT). A solver written
for a different dataset, carried across without a line changed, beats all three.
It does not approach the native solver's $64.54\%$, and it would not have
attained $\max\mathcal{B}$ in any reported group---but it establishes that the
control-group requirement does not depend on bespoke engineering for the exact
benchmark under test. Half of a symbolic baseline is portable; the table parser
is not.
